% บทคัดย่อภาษาไทย
% ตัวเลขความคืบหน้าปรับปรุงให้ตรงกับบทที่ 5 ฉบับปัจจุบัน (สถานะ ณ เดือนสิงหาคม 2569)
% ส่วนหัว (ชื่อผู้จัดทำ/ปี/ที่ปรึกษา) และคำสำคัญท้ายบทคัดย่อ สร้างอัตโนมัติจาก main.tex ไม่ต้องใส่ในไฟล์นี้

โครงงานนี้มีวัตถุประสงค์เพื่อพัฒนาเว็บแอปพลิเคชัน ``ไดโน'' (Dino) ให้เป็นแพลตฟอร์มศูนย์กลางข้อมูลการท่องเที่ยวในจังหวัดขอนแก่น เพื่อแก้ไขปัญหาการเข้าถึงข้อมูลการท่องเที่ยวที่กระจัดกระจายและไม่ครบถ้วน โครงงานนี้ได้บูรณาการเทคโนโลยีปัญญาประดิษฐ์ (AI) ผ่านโมเดลภาษาขนาดใหญ่ (LLMs) ร่วมกับสถาปัตยกรรม Retrieval-Augmented Generation (RAG) ในการพัฒนาระบบหลัก 2 ส่วน ได้แก่ (1) ระบบสร้างแผนการเดินทางแบบปรับตัว (Adaptive Trip Planning) ที่ประเมินจากข้อจำกัดด้านงบประมาณ ระยะเวลา และความสนใจเฉพาะบุคคล พร้อมกลไกตรวจสอบคุณภาพแผนด้วยเทคนิค LLM-as-a-Judge และ (2) ระบบแชตบอต (Chatbot) สำหรับโต้ตอบและให้คำแนะนำด้านการท่องเที่ยวแบบเรียลไทม์ โดยใช้ฐานข้อมูลสถานที่และร้านอาหารจาก Google Places API ภายในจังหวัดขอนแก่น

ผลการดำเนินงาน ณ ปัจจุบันพบว่า ระบบหลังบ้าน ระบบส่วนหน้า และบริการปัญญาประดิษฐ์เชื่อมต่อและทำงานร่วมกันได้จริงแบบครบวงจร (End-to-End) แล้ว ทั้งฟังก์ชันแชทบอทและระบบวางแผนการเดินทางด้วย LLM ซึ่งมีความคืบหน้ารวมประมาณร้อยละ 85--90 ของขอบเขตที่ออกแบบไว้ อย่างไรก็ตาม ระบบยืนยันตัวตนและระบบสแกน QR สะสม/แลกพอยท์ยังอยู่ในลักษณะต้นแบบ (Prototype) ที่ยังไม่เชื่อมต่อกับฐานข้อมูลจริง และยังไม่ได้ผ่านการทดสอบกับผู้ใช้งานจริงในวงกว้าง (User Acceptance Testing) ทั้งนี้ เมื่อการพัฒนาส่วนที่เหลือและการทดสอบกับผู้ใช้งานจริงเสร็จสมบูรณ์ คาดว่าแอปพลิเคชันนี้จะช่วยลดระยะเวลาในการค้นหาข้อมูล เพิ่มประสิทธิภาพในการวางแผนการเดินทางที่ตอบโจทย์เฉพาะบุคคล ตลอดจนช่วยสนับสนุนผู้ประกอบการท้องถิ่นและส่งเสริมเศรษฐกิจการท่องเที่ยวของจังหวัดขอนแก่นอย่างยั่งยืน
